\section{Background}\label{sec:background}

A robot that must know where it is has two unreliable sources of information: the controls it was given, which are
corrupted by wheel slip, and the measurements it takes of its surroundings, which are corrupted by sensor noise.
Combining them into an estimate of the pose, with an honest account of its uncertainty, is probabilistic
localization; doing it while building the map is SLAM \cite{thrun2005probabilistic}. This section gives what is
needed to follow the rest of the report: the network the problem is written on, how inference is done on it
exactly and how it has to be approximated, what makes that approximation affordable in SLAM, and finally what it
means to learn the model's numbers rather than supply them.

\subsection{Localization and SLAM as a dynamic Bayesian network}\label{sec:dbn}

In localization the map is known and the robot estimates its pose from the controls it was given and its
observations of known landmarks. In SLAM the map is unknown, and the robot estimates the map and its own trajectory at the same time.

Both are naturally written as a dynamic Bayesian network \cite[ch.~6 and 15]{koller2009probabilistic}. The pose
$\x_t = (x_t, y_t, \theta_t)$ depends on the previous pose and the control $\uu_t$ through the motion model
$p(\x_t \mid \x_{t-1}, \uu_t)$; each observation $\z_t$ is conditioned on the pose from which it was taken and on
the map $\m$ through the measurement model $p(\z_t \mid \x_t, \m)$ (\cref{fig:dbn}), so
\begin{equation}
  p(\x_{0:T}, \m, \z_{1:T} \mid \uu_{1:T}) = p(\x_0)\, p(\m) \prod_{t=1}^{T} p(\x_t \mid \x_{t-1}, \uu_t)\,
  p(\z_t \mid \x_t, \m).
  \label{eq:joint}
\end{equation}
Controls and observations are observed; the trajectory is latent, and in SLAM the map is latent as well.
Correspondences are known throughout, so each observation carries the identity of the landmark it measures.

\begin{figure}[!ht]
  \centering
  \setlength{\unitlength}{1mm}
  \begin{picture}(110,42)
    \put(10,20){\circle{10}}\put(10,20){\makebox(0,0){$\x_{t-2}$}}
    \put(40,20){\circle{10}}\put(40,20){\makebox(0,0){$\x_{t-1}$}}
    \put(70,20){\circle{10}}\put(70,20){\makebox(0,0){$\x_t$}}
    \put(15,20){\vector(1,0){20}}\put(45,20){\vector(1,0){20}}
    \put(40,37){\circle{8}}\put(40,37){\makebox(0,0){$\uu_{t-1}$}}
    \put(70,37){\circle{8}}\put(70,37){\makebox(0,0){$\uu_t$}}
    \put(40,33){\vector(0,-1){8}}\put(70,33){\vector(0,-1){8}}
    \put(40,3){\circle{8}}\put(40,3){\makebox(0,0){$\z_{t-1}$}}
    \put(70,3){\circle{8}}\put(70,3){\makebox(0,0){$\z_t$}}
    \put(40,15){\vector(0,-1){8}}\put(70,15){\vector(0,-1){8}}
    \put(100,20){\circle{10}}\put(100,20){\makebox(0,0){$\m$}}
    \put(95.5,18){\vector(-2,-1){21.5}}\put(95.5,17.5){\vector(-4,-1){52}}
  \end{picture}
  \caption{Two time slices of the network. Controls $\uu$ and observations $\z$ are observed; poses $\x$ and the
    map $\m$ are latent. In localization the map is known.}
  \label{fig:dbn}
\end{figure}

\subsection{Exact inference on the network: the Bayes filter}\label{sec:bayesfilter}

Estimation on this network is the Bayes filter, the forward pass of inference on the network written out over
time. What it computes is the \emph{belief}, the distribution over the current pose given everything observed so
far. Because $\x_t$ separates past from future, the belief satisfies
\begin{equation}
  p(\x_t \mid \z_{1:t}, \uu_{1:t}) \;\propto\; p(\z_t \mid \x_t)
  \int p(\x_t \mid \x_{t-1}, \uu_t)\, p(\x_{t-1} \mid \z_{1:t-1}, \uu_{1:t-1})\, d\x_{t-1},
  \label{eq:bayes-filter}
\end{equation}
a prediction through the motion model followed by a correction with the new observation
\cite{thrun2005probabilistic}.

The state is continuous, so the belief is a density and the prediction is an integral. That integral is tractable
in two cases: with a discrete state space it is a finite sum, and \eqref{eq:bayes-filter} is the forward algorithm
of a hidden Markov model; with a linear model and Gaussian noise the belief stays Gaussian and the recursion is
the Kalman filter. Range and bearing measurements are nonlinear in the pose, and under global uncertainty the
belief is multimodal, so neither applies and no closed form exists.

\subsection{Approximate inference: the particle filter}\label{sec:pf}

Where the integral cannot be carried out it is replaced by samples. The distribution the samples are drawn from is
the \emph{proposal}. Drawing a pose from the motion model and weighting it by $p(\z_t \mid \x_t)$ is likelihood
weighting applied to one slice, with the prior as the proposal \cite[ch.~12]{koller2009probabilistic}; carrying
the weighted set forward to the next slice makes it sequential \cite{gordon1993novel,arulampalam2002tutorial}. The belief after step $t$ is $N$ particles
$\x_t^{(i)}$ with normalized weights $\w_t^{(i)}$.

One ingredient has to be added. Over time the weights become products of more and more terms, so
they concentrate on a single particle and the set stops representing the belief; how far this has gone is measured
by the effective sample size $\mathrm{ESS} = 1/\sum_i (\w^{(i)})^2$. Resampling, which replaces the set by $N$
draws from itself, prevents it, and has no counterpart in the exact algorithm. It costs something in return:
drawing with replacement duplicates the heaviest particles and discards the rest, so the number of distinct
particles falls. A filter therefore has four things to decide: the proposal, when to resample, how to
resample, and how many particles to carry. The middle two are the subject of this report.

\clearpage
\subsection{Rao-Blackwellization and FastSLAM}\label{sec:rbpf}

Applying the particle filter to SLAM raises a problem of dimension. The state is now the pose together with the
map, so each particle would have to carry a complete map as well as a pose, and the state space gains two
dimensions per landmark. Sampling that joint space directly is intractable: the number of particles needed to
cover it grows exponentially with the number of landmarks.

The structure of the network offers a way out. An observation of landmark $\m_k$ has exactly two parents: the
pose it was taken from and $\m_k$ itself. Any path between two landmarks must therefore leave one of them
through an observation and pass through that observation's pose, with both edges pointing away from the pose, and
a known trajectory blocks every such path. The landmarks are d-separated given the trajectory and the
observations \cite{koller2009probabilistic}, so the posterior splits:
\begin{equation}
  p(\x_{1:t}, \m \mid \z_{1:t}, \uu_{1:t}) = p(\x_{1:t} \mid \z_{1:t}, \uu_{1:t})
  \prod_{k=1}^{K} p(\m_k \mid \x_{1:t}, \z_{1:t}).
  \label{eq:rao-blackwell}
\end{equation}
The first factor is a distribution over trajectories, the hard part. Each of the others is a small problem over
one landmark's two coordinates, given that trajectory. Splitting a posterior in this way and solving the easy
part exactly, rather than by sampling it, is Rao-Blackwellization \cite{doucet2000rao}, and FastSLAM is its
application here: every particle carries one trajectory and one extended Kalman filter per landmark
\cite{montemerlo2002fastslam}. Only the trajectory is sampled, which is what makes particle-filter SLAM
affordable.

\subsection{Learning the parameters of the network}\label{sec:learning}

Everything above takes the model as fully specified and asks how to compute with it. A graphical model raises a
second question: where its numbers come from.

The structure of \eqref{eq:joint} fixes which distributions appear, not the numbers inside them. Write those
numbers as $\theta$. They may be supplied, or estimated from data; estimating them is \emph{parameter} learning,
as opposed to structure learning, which asks for the graph itself.

When some of the variables are never observed it is parameter learning \emph{with incomplete data}
\cite[ch.~19]{koller2009probabilistic}, and that is the case here: the trajectory the observations depend on is
latent, so the only data are the controls and the observations. A dynamic Bayesian network is also a template
model, so the same conditional distribution is instantiated at every time slice and its parameters are tied
across slices \cite[ch.~6]{koller2009probabilistic}. Every step is then evidence about the same $\theta$, and a
longer trajectory is more data about it rather than more unknowns.

The estimate is the $\theta$ under which the observations are most probable,
\begin{equation}
  \hat{\theta} = \arg\max_{\theta}\; \log p(\z_{1:T} \mid \uu_{1:T}, \theta), \qquad
  p(\z_{1:T} \mid \uu_{1:T}, \theta) = \prod_{t=1}^{T} p(\z_t \mid \z_{1:t-1}, \uu_{1:t}, \theta).
  \label{eq:mle}
\end{equation}
Its maximum is in the right place and not merely a convenient one: the average log likelihood tends to a term
free of $\theta$ minus the Kullback-Leibler divergence between the true distribution and the model, and that
divergence vanishes only where the two agree \cite[ch.~17]{koller2009probabilistic}.

Each factor of \eqref{eq:mle} is again an integral with no closed form, and the particle filter estimates it as a
by-product. With particles propagated through the motion model the unnormalized weight is $\tilde{\w}_t^{(i)} =
\w_{t-1}^{(i)}\, p(\z_t \mid \x_t^{(i)}, \theta)$, and
\begin{equation}
  p(\z_t \mid \z_{1:t-1}, \uu_{1:t}, \theta) \;\approx\; \sum_{i=1}^{N} \tilde{\w}_t^{(i)},
  \label{eq:evidence}
\end{equation}
the sum the filter forms anyway in order to normalize \cite{doucet2011tutorial}.

Two consequences follow, and both appear in the results. The estimate of the likelihood is unbiased, so the
estimate of its logarithm is biased downwards, and the bias grows as the particle approximation degrades: the
quality of the inference limits what can be learned. And the curvature of the expected score at $\theta^*$ is
$T\,\mathcal{I}(\theta^*)$, with $\mathcal{I}$ the information one step carries, so the variance of the estimate
is $1/(T\mathcal{I})$. A parameter the observations say little about leaves a flat likelihood: it is in the model
and still not identifiable from the data, and because the parameters are tied, a longer trajectory sharpens it.
